% ======================================================================
\section{States, normalisation and matrix elements}\label{app:states}
% ======================================================================

\subsection{States and normalisation}

The soft vacuum $\ket0$ is the state annihilated by all soft annihilation operators, $a^a_i(k^+,\v k)\ket0=0$ for $\Lambda<k^+<\vee$. It is a state of the soft modes only. The complete state of the projectile is a vector in the product of the soft and the valence Hilbert spaces, $\mathcal H_{\rm soft}\otimes\mathcal H_{\rm val}$, and $\ket0$ is to be read as $\ket0\otimes\ket v$ with $\ket v$ the state of the valence modes, which is left unspecified. The soft vacuum is therefore not the vacuum of the theory. Its energy is set to zero, $E_0=0$. The one- and two-gluon states are
\begin{align}\label{eq:states}
\ket{g^a_i(k)}&=\frac{a^{\dagger a}_i(k)}{(2\pi)^{3/2}}\ket0\,,&
\ket{g^a_i(k)\,g^b_j(p)}&=\frac{a^{\dagger a}_i(k)\,a^{\dagger b}_j(p)}{(2\pi)^3}\ket0\,,\notag\\
\ket{g^a_i(k^+,\v z)}&=\frac{a^{\dagger a}_i(k^+,\v z)}{(2\pi)^{1/2}}\ket0\,,&
\ket{g^a_i(k^+,\v z)\,g^b_j(p^+,\v z')}&=\frac{a^{\dagger a}_i(k^+,\v z)\,a^{\dagger b}_j(p^+,\v z')}{2\pi}\ket0\,,
\end{align}
and states with more gluons are defined in the same way, with a factor $(2\pi)^{-3/2}$ for each gluon in momentum space. They are eigenstates of $H_0$ with the energies $E_g(k)$ and $E_{gg}(k,p)$ of Eq.~\eqref{eq:measure}, and they are normalised to $\delta$-functions without factors of $2\pi$,
\begin{align}\label{eq:norms}
\braket{g^b_j(p)|g^a_i(k)}&=\delta^{ab}\delta_{ij}\,\delta^{(3)}(k-p)\,,\notag\\
\braket{g^b_j(p^+,\v z')|g^a_i(k^+,\v z)}&=\delta^{ab}\delta_{ij}\,\delta(k^+-p^+)\,\delta^{(2)}(\v z-\v z')\,,\notag\\
\braket{g^{a'}_{i'}(k')\,g^{b'}_{j'}(p')|g^a_i(k)\,g^b_j(p)}&=\delta^{aa'}\delta_{ii'}\delta^{(3)}(k-k')\,\delta^{bb'}\delta_{jj'}\delta^{(3)}(p-p')\notag\\
&\quad+\delta^{ab'}\delta_{ij'}\delta^{(3)}(k-p')\,\delta^{ba'}\delta_{ji'}\delta^{(3)}(p-k')\,.
\end{align}
The states in transverse coordinate space are related to those in momentum space by
\begin{align}\label{eq:stateFT}
\ket{g^a_i(k^+,\v z)}=\int\frac{d^2k}{2\pi}\,e^{-i\v k\cdot\v z}\ket{g^a_i(k^+,\v k)}\,,\qquad
\ket{g^a_i(k^+,\v k)}=\int\frac{d^2z}{2\pi}\,e^{i\v k\cdot\v z}\ket{g^a_i(k^+,\v z)}\,.
\end{align}
A two-gluon state is symmetric under the exchange of the labels of the two gluons. An expansion $\int d^3p\,d^3q\,F^{bc}_{jk}(p,q)\ket{g^b_j(p)\,g^c_k(q)}$ therefore determines only the part of $F$ that is symmetric under $(b,j,p)\leftrightarrow(c,k,q)$, and the sum over the two-gluon states with unit weight is $\frac12\int d^3p\,d^3q\ket{g^b_j(p)g^c_k(q)}\bra{g^b_j(p)g^c_k(q)}$.


\subsection{Matrix elements}\label{app:matrixelements}

With the states~\eqref{eq:states}, the matrix elements needed in this paper are
\begin{align}
\braket{g^a_i(k)|H_g|0}&=\frac{g\,k^i\rho^a(-\v k)}{4\pi^{3/2}(k^+)^{3/2}}\,,\label{eq:m1}\\
\braket{g^c_l(q)\,g^b_j(p)|H_g|g^a_i(k)}&=\delta^{ac}\delta_{il}\,\delta^{(3)}(k-q)\,\frac{g\,p^j\rho^b(-\v p)}{4\pi^{3/2}(p^+)^{3/2}}
+\delta^{ab}\delta_{ij}\,\delta^{(3)}(k-p)\,\frac{g\,q^l\rho^c(-\v q)}{4\pi^{3/2}(q^+)^{3/2}}\,,\label{eq:m2}\\
\braket{g^c_l(q)\,g^b_j(p)|H_{ggg}|g^a_i(k)}&=\frac{igf^{abc}\,\delta^{(3)}(k-p-q)}{8\pi^{3/2}\sqrt{k^+p^+q^+}}\Bigg[\Big(p^i-q^i+\frac{q^+-p^+}{k^+}k^i\Big)\delta_{jl}\notag\\
&\hspace{0.6cm}+\Big(k^j+q^j-\frac{k^++q^+}{p^+}p^j\Big)\delta_{il}
+\Big(\frac{k^++p^+}{q^+}q^l-p^l-k^l\Big)\delta_{ij}\Bigg]\,,\label{eq:m3}\\
\braket{g^c_j(q)\,g^b_i(p)|H_{gg\text{-inst}}|0}&
=\frac{ig^2f^{abc}\,(q^+-p^+)\,\delta_{ij}\,\rho^a(-\v p-\v q)}{2(2\pi)^3\sqrt{p^+q^+}\,(p^++q^+)^2}\,,\label{eq:m4}\\
\braket{g^c_j(q)|H_{gg\text{-inst}}|g^b_i(p)}&
=\frac{ig^2f^{abc}\,(p^++q^+)\,\delta_{ij}\,\rho^a(\v p-\v q)}{2(2\pi)^3\sqrt{p^+q^+}\,(p^+-q^+)^2}\,.\label{eq:m5}
\end{align}
Eq.~\eqref{eq:m1} is the emission of a soft gluon by the valence charge; Eq.~\eqref{eq:m2} the emission of a second gluon in the presence of a first one, which is a spectator; Eq.~\eqref{eq:m3} the splitting of a soft gluon into two; Eq.~\eqref{eq:m4} the instantaneous emission of two gluons; Eq.~\eqref{eq:m5} the instantaneous interaction of a soft gluon with the valence charge. The matrix elements of the opposite transitions follow by Hermitian conjugation, with $[\rho^a(\v p)]^\dagger=\rho^a(-\v p)$; for instance $\braket{0|H_g|g^a_i(k)}=g\,k^i\rho^a(\v k)/[4\pi^{3/2}(k^+)^{3/2}]$. Eqs.~\eqref{eq:m1}, \eqref{eq:m2}, \eqref{eq:m4} and~\eqref{eq:m5} follow directly from Eqs.~\eqref{eq:Hg} and~\eqref{eq:Hinst}; all of them are operators on $\mathcal H_{\rm val}$ except Eq.~\eqref{eq:m3}.

As an illustration, Eq.~\eqref{eq:m1} follows from Eq.~\eqref{eq:Hg} with $\bra{g^a_i(k)}a^{\dagger b}_j(p)\ket0=(2\pi)^{3/2}\delta^{ab}\delta_{ij}\delta^{(3)}(k-p)$: the factor is $(2\pi)^{3/2}/(2\pi)^3\times g/\sqrt2=g/(4\pi^{3/2})$. In Eq.~\eqref{eq:m2} the gluon that is not emitted is a spectator, and the two terms correspond to the two ways of identifying it with one of the outgoing gluons. Eqs.~\eqref{eq:m4} and~\eqref{eq:m5} follow in the same way from the first and the last term of Eq.~\eqref{eq:Hinst}; in Eq.~\eqref{eq:m4} the two contractions of the pair of creation operators combine into the factor $q^+-p^+$ by the antisymmetry of $f^{abc}$. Eq.~\eqref{eq:m3} is taken from~\cite{Lublinsky:2016meo}.
